% pdflatex -shell-escape fig_01.tex
% =============================================================================
\documentclass[multi,crop,landscape]{standalone}
\standaloneconfig{border=0pt}

\usepackage{amsmath} % needed for \dfrac
\usepackage{pgfplots}
\usepgfplotslibrary{groupplots}
\pgfplotsset{compat=1.18, lua backend=false}

% ---------- Safe helper ----------
\pgfmathdeclarefunction{clipone}{1}{%
  \pgfmathparse{min(max(#1,0.0001),0.9999)}%
}

% ---------- P(u) functions ----------
% Uniform(0,1): P(u) = (1+u)/(2u)
\pgfmathdeclarefunction{Puni}{1}{%
  \pgfmathparse{ (1 + clipone(#1)) / (2 * clipone(#1)) }%
}
% -------------- b > a -----

% P(u) for Uniform(a,b), assuming b > a and a,b > 0
% Usage: PuniAB(u, a, b)
\pgfmathdeclarefunction{PuniAB}{3}{%
  \pgfmathparse{
    ( #2 + ( (#3 - #2) * (1 + clipone(#1)) ) / 2 ) /
    ( #2 + ( (#3 - #2) * clipone(#1) ) )
  }%
}
% -------------------------

% Exponential: P(u) = 1 - 1 / ln(1-u)
\pgfmathdeclarefunction{Pexp}{1}{%
  \pgfmathparse{ 1 - 1 / ln(1 - clipone(#1)) }%
}

% Power: P(u) = [ beta(1 - u^{1+1/beta}) ] / [ (1-u) u^{1/beta} (beta+1) ]
\pgfmathdeclarefunction{Ppow}{2}{%
  \pgfmathparse{ (#2*(1 - clipone(#1)^(1 + 1/#2))) /
                 ( (1 - clipone(#1)) * clipone(#1)^(1/#2) * (#2 + 1) ) }%
}

\begin{document}
\begin{tikzpicture}
  \begin{groupplot}[
    group style={group size=2 by 2, horizontal sep=2.0cm, vertical sep=3cm},
    width=12.5cm, height=9.5cm,
    samples=2000,
    grid=both, major grid style={gray!25}, minor grid style={gray!10},
    xlabel={$u$}, ylabel={$P(u)$},
    title style={font=\small, align=center},
    legend style={font=\small, draw=none, fill=none},
    every axis/.append style={thick}
  ]

  % ================= Panel 1: Uniform (fixed a=0, b=1) =================
  \nextgroupplot[
    title={Uniform $(a{=}0,b{=}1)$:\quad $P(u)=\dfrac{1+u}{2u}$},
    % xmin=0.1, xmax=0.99, domain=0.1:0.99
    xmin=0, xmax=1, ymin=0, ymax=10,
    restrict y to domain=0:10,
    unbounded coords=jump
  ]
    % \addplot[black, dashed, line width=2pt] {1};
    % \addplot[blue, line width=2pt] { Puni(x) }; \addlegendentry{Uniform $(0,1)$ (mean $=0.5$)}
    
    \addplot[black, dashed, line width=2pt] {1}; \addlegendentry{reference line at 1}
    \addplot[blue, line width=2pt] { Puni(x) }; \addlegendentry{Uniform $(0,1)$}
    \addplot[red, line width=2pt] { PuniAB(x,1,2) }; \addlegendentry{Uniform $(1,2)$}
    \addplot[teal!70!black, line width=2pt] { PuniAB(x,1,3) }; \addlegendentry{Uniform $(1,3)$}

  % ================= Panel 2: Exponential =================
  % Means 1,2,3 with ?? = 1, 0.5, 1/3 (P(u) does not depend on ??).
  \nextgroupplot[
    title={Exponential $(\lambda)$:\quad $P(u)=1-\dfrac{1}{\log(1-u)}$},
    % xmin=0.2, xmax=0.8, domain=0.2:0.8
    % xmin=0.1, xmax=0.99, domain=0.1:0.99
    xmin=0, xmax=1, ymin=0, ymax=10,
    restrict y to domain=0:10,
    unbounded coords=jump
  ]
    \addplot[black, dashed, line width=2pt] {1}; \addlegendentry{reference line at 1}
    \addplot[blue, line width=2pt] { Pexp(x) }; \addlegendentry{\ $\lambda=1$}
    % \addplot[red, line width=2pt]  { Pexp(x) }; \addlegendentry{mean=2;\ $\lambda=0.5$}
    % \addplot[teal!70!black, line width=2pt] { Pexp(x) }; \addlegendentry{mean=3;\ $\lambda=\frac{1}{3}$}

  % ================= Panel 3: Pareto I =================
  % ??=1. Mean = ????/(??-1) => feasible for mean>1 only.
  \nextgroupplot[
    title={Pareto I $(\beta{=}1,\alpha)$:\quad $P(u)=\dfrac{\alpha}{\alpha-1}$},
    % xmin=0.1, xmax=0.99, domain=0.05:0.99
    % xmin=0.1, xmax=0.99, domain=0.1:0.99
    xmin=0, xmax=1, ymin=0, ymax=10,
    restrict y to domain=0:10,
    unbounded coords=jump
  ]
    \addplot[black, dashed, line width=2pt] {1}; \addlegendentry{reference:\ $P{=}1$} % <--- add legend entry
    \addplot[blue, line width=2pt] { 2.0 / (2.0 - 1.0) }; \addlegendentry{\ $\alpha=2$}
    \addplot[red, line width=2pt]  { 1.5 / (1.5 - 1.0) }; \addlegendentry{\ $\alpha=1.5$}
    % (mean=1 not possible for Pareto I)

  % ================= Panel 4: Power =================
  % Mean = a??/(??+1). Fix ??=1 and set a = 2,4,6 for means 1,2,3.
  % P(u) depends on ?? (not on a), so curves coincide for ??=1.
  \nextgroupplot[
    title={Power $(a,\beta)$:\quad $P(u)=\dfrac{\beta\,(1-u^{1+1/\beta})}{(1-u)\,u^{1/\beta}\,(\beta+1)}$},
    % xmin=0.1, xmax=0.99, domain=0.1:0.99
    xmin=0, xmax=1, ymin=0, ymax=10,
    restrict y to domain=0:10,
    unbounded coords=jump
  ]
    \addplot[black, dashed, line width=2pt] {1}; \addlegendentry{reference line at 1}
    \addplot[blue, line width=2pt] { Ppow(x,1.0) }; \addlegendentry{$\beta=1$}
    \addplot[red, line width=2pt]  { Ppow(x,2.0) }; \addlegendentry{$\beta=2$}
    \addplot[teal!70!black, line width=2pt] { Ppow(x,3.0) }; \addlegendentry{$\beta=3$}

  \end{groupplot}
\end{tikzpicture}
\end{document}

% =============================================================================

% --- Post-build conversion to PNG via ImageMagick ---
% Transparent PNG (recommended for slides):
% \AtEndDocument{%
%   \immediate\write18{magick convert -density 300 -background none -alpha on "\fig_01.pdf" "\fig_01.png"}%
% }
% If you prefer opaque (white) background, use instead:
\AtEndDocument{%
  \immediate\write18{magick convert -density 300 "\fig_01.pdf" "\fig_01.png"}%
}

% =============================================================================
